\section{Resultat numérique}

\subsection{Schéma RKDG et reconstruit}

Si l'on essaye de comparer les deux schéma, d'avoir les mêmes conditions et donc la même initialisation.
On comparera les deux représentation sur divers problème classique de la littérature. 
\subsubsection{advection \red{à changer}}
Le problème de l'advection défini par $f(u)= c\, u$ a l'avantage d'être linéaire et correspond à un déplacement d'un signal à vitesse $c$.
La solution exacte est on ne peut plus simple à calculer comme $u_{ex}(x,t) =u_0(x-ct)$. \\
Ces deux charactéristiques rendent ce problème à la fois simple à étudier et très utile pour tester la convergence d'un modèle dans des cas simples.
\paragraph{test régulier : signal sinusoidale}
$u_0(x) = sin(2\pi x)$
\paragraph{test de dispersion : signal composite}
On regardera ici le cas d'un signal composite premièrement introduit dans \cite{JIANG1996202}, 
est une solution initiale composé de 4 signaux distinct permettant de tester la \red{dispersion} du schéma pour différents signaux. 
Une combinaison de gaussiène, un signal carré, un signal triangulaire et une combinaison de racine carré.
\begin{equation}
    u_0(x) = \left\{
\begin{aligned}
    &\frac{1}{6}(G(x,\beta,z-\delta) +G(x,\beta,z+\delta)+4G(x,\beta,z)) & \text{si } -0.8 \leq x\leq -0.6\\
    &1 & \text{si } -0.4 \leq x \leq -0.2 \\ 
    &1-|10(x-0.1)| & \text{si } 0 \leq x \leq 0.2 \\
    &\frac{1}{6}(F(x,\alpha,a-\delta) +F(x,\alpha,a+\delta)+4F(x,\alpha,a)) & \text{si } 0.4\leq x \leq 0.6 \\
    &0 & \text{sinon }
\end{aligned}
    \right.
\end{equation}

\subsubsection{Burgers \red{à changer}}
Le problème de Burgers défini par $f(u) = \frac{u^2}{2}$ est l'un des plus simple problème non linéaire, 
et étant donnée une solution régulière simple comme une signal sinusoidale $u_0(x)= \sin(2\pi x)$ ou un problème de Riemann,
La solution exacte est relativement simple à trouver. 
Dans le cas du signal sinusoidale, on utilisera une méthode de dichotomie \red{explication ?} pour retrouver la solution, cette méthode ne fonctionnant uniquement pour $t$ inférieur à un temps \textit{critique}.\\
Cela nous donnera un cas test pour un flux non linéaire simple.
\paragraph{cas test régulier : sinusoidale}

\paragraph{cas test simple : combinaison de problème de Riemann}

\subsubsection{Buckley \red{à changer}}
les deux cas test étant introduit dans \cite{CHEN2017427}
\paragraph{cas test Riemann}
\paragraph{cas test carré}

\subsubsection{Système \red{à changer}}

\subsection{Schéma Monolithique}

\subsubsection{GMP/LMP \red{à changer}}
\subsubsection{positivité \red{à changer}}
\subsubsection{entropie \red{à changer}}